% ============================================================
%  BioOS Causal Constitution — Bilingual Scientific Paper
%  Szőke László-Ferenc · Citrom Média LTD · 2026-05-24
%  Engine: XeLaTeX
% ============================================================
\documentclass[11pt,a4paper,twoside]{article}

% ── XeLaTeX fonts ────────────────────────────────────────────
\usepackage{fontspec}
\defaultfontfeatures{Ligatures=TeX}
\setmainfont{Latin Modern Roman}
\setmonofont[Scale=MatchLowercase]{Latin Modern Mono}
\setsansfont{Latin Modern Sans}

% ── Language ─────────────────────────────────────────────────
\usepackage{polyglossia}
\setmainlanguage{english}
\setotherlanguage{magyar}

% ── Math ─────────────────────────────────────────────────────
\usepackage{amsmath,amsthm,amssymb,mathtools}

% ── Layout ───────────────────────────────────────────────────
\usepackage[top=2.5cm,bottom=2.5cm,left=3cm,right=3cm,headheight=14pt]{geometry}
\usepackage{fancyhdr}
\usepackage{titlesec}
\usepackage{float}
\usepackage{caption}
\usepackage{enumitem}

% ── Colors & Boxes ───────────────────────────────────────────
\usepackage{xcolor}
\usepackage{mdframed}

% ── Tables ───────────────────────────────────────────────────
\usepackage{booktabs}
\usepackage{array}
\usepackage{longtable}
\usepackage{tabularx}
\usepackage{multirow}

% ── Code listings ────────────────────────────────────────────
\usepackage{listings}

% ── Graphics ─────────────────────────────────────────────────
\usepackage{tikz}
\usetikzlibrary{arrows.meta,positioning,shapes.geometric,fit,backgrounds,automata}

% ── Hyperlinks ───────────────────────────────────────────────
\usepackage[
  colorlinks=true,
  linkcolor=bioblue,
  urlcolor=biogreen,
  citecolor=biogreen,
  unicode=true,
  pdftitle={BioOS Causal Constitution},
  pdfauthor={Szoke Laszlo-Ferenc},
]{hyperref}

% =============================================================
%  Colours
% =============================================================
\definecolor{bioblue}{RGB}{0,70,140}
\definecolor{biogreen}{RGB}{27,140,80}
\definecolor{biored}{RGB}{170,35,25}
\definecolor{biogray}{RGB}{80,85,95}
\definecolor{codebg}{RGB}{22,27,34}
\definecolor{codetext}{RGB}{201,209,217}
\definecolor{codecomment}{RGB}{130,140,150}
\definecolor{codestring}{RGB}{121,192,255}
\definecolor{codekw}{RGB}{255,123,114}
\definecolor{codekw2}{RGB}{210,168,255}
\definecolor{codenumber}{RGB}{240,198,116}
\definecolor{lightbluebg}{RGB}{235,242,250}
\definecolor{lightgreenbg}{RGB}{232,248,238}
\definecolor{lightyellowbg}{RGB}{255,252,230}

% =============================================================
%  Macros
% =============================================================
\newcommand{\sem}[1]{\mathopen{[\![}#1\mathclose{]\!]}}
\newcommand{\nt}[1]{\textit{#1}}
\newcommand{\tm}[1]{\texttt{#1}}
\newcommand{\bnfor}{\;\mid\;}
\newcommand{\bnfdef}{\;{::=}\;}
\newcommand{\bioURL}{\url{https://bioos.metaspace.bio}}
\newcommand{\proofURL}{\url{https://bioos.metaspace.bio/proof}}
\newcommand{\irule}[3]{\dfrac{\displaystyle #1}{\displaystyle #2}\quad\textsc{(#3)}}

% =============================================================
%  Theorem environments
% =============================================================
\mdfdefinestyle{thmbox}{
  linewidth=1.5pt, linecolor=bioblue, backgroundcolor=lightbluebg,
  innertopmargin=9pt, innerbottommargin=9pt,
  innerleftmargin=12pt, innerrightmargin=12pt,
  skipabove=8pt, skipbelow=6pt,
}
\mdfdefinestyle{defbox}{
  linewidth=1pt, linecolor=biogreen, backgroundcolor=lightgreenbg,
  innertopmargin=8pt, innerbottommargin=8pt,
  innerleftmargin=12pt, innerrightmargin=12pt,
  skipabove=8pt, skipbelow=6pt,
}
\mdfdefinestyle{notebox}{
  linewidth=1pt, linecolor=biogray, backgroundcolor=lightyellowbg,
  innertopmargin=7pt, innerbottommargin=7pt,
  innerleftmargin=10pt, innerrightmargin=10pt,
  skipabove=6pt, skipbelow=4pt,
}
% English
\newmdtheoremenv[style=thmbox]{theorem}{Theorem}[section]
\newmdtheoremenv[style=thmbox]{corollary}[theorem]{Corollary}
\newmdtheoremenv[style=defbox]{definition}[theorem]{Definition}
\newtheorem*{remark}{Remark}
% Hungarian
\newmdtheoremenv[style=thmbox]{tetel}{Tétel}[section]
\newmdtheoremenv[style=thmbox]{kovetkezmeny}[tetel]{Következmény}
\newmdtheoremenv[style=defbox]{definicio}[tetel]{Definíció}
\newtheorem*{megjegyzes}{Megjegyzés}

% =============================================================
%  Listings
% =============================================================
\lstdefinelanguage{bio}{
  keywords={CELL,INTERFACE,INVARIANTS,STATES,STATE,RULE,INPUT,OUTPUT,
            TRANSITION,TO,IF,TYPE,CRITICAL\_SHIELD},
  keywords=[2]{AND,OR,NOT,IMPLIES,TRUE,FALSE},
  keywords=[3]{BINARY,INTEGER,FLOAT,TIMESTAMP,VECTOR2D,VECTOR3D,ENUM},
  sensitive=true,
  morecomment=[l]{//},
  morestring=[b]",
}
\lstdefinestyle{biostyle}{
  language=bio,
  backgroundcolor=\color{codebg},
  basicstyle=\ttfamily\footnotesize\color{codetext},
  keywordstyle=\color{codekw}\bfseries,
  keywordstyle=[2]\color{codekw2}\bfseries,
  keywordstyle=[3]\color{codenumber},
  commentstyle=\color{codecomment}\itshape,
  stringstyle=\color{codestring},
  breaklines=true, frame=single, rulecolor=\color{gray!50},
  xleftmargin=0.5em, xrightmargin=0.5em, framesep=5pt,
  numbers=left, numberstyle=\tiny\color{codecomment}, numbersep=7pt,
  tabsize=2, showstringspaces=false, aboveskip=8pt, belowskip=8pt,
}
\lstdefinestyle{shellstyle}{
  backgroundcolor=\color{codebg},
  basicstyle=\ttfamily\footnotesize\color{codetext},
  breaklines=true, frame=single, rulecolor=\color{gray!50},
  xleftmargin=0.5em, xrightmargin=0.5em, framesep=5pt,
  showstringspaces=false, aboveskip=6pt, belowskip=6pt,
}
\lstdefinestyle{pythonstyle}{
  language=Python,
  backgroundcolor=\color{codebg},
  basicstyle=\ttfamily\footnotesize\color{codetext},
  keywordstyle=\color{codekw}\bfseries,
  commentstyle=\color{codecomment}\itshape,
  stringstyle=\color{codestring},
  breaklines=true, frame=single, rulecolor=\color{gray!50},
  xleftmargin=0.5em, xrightmargin=0.5em, framesep=5pt,
  numbers=left, numberstyle=\tiny\color{codecomment}, numbersep=7pt,
  tabsize=4, showstringspaces=false, aboveskip=6pt, belowskip=6pt,
}

% =============================================================
%  Page style
% =============================================================
\pagestyle{fancy}
\fancyhf{}
\fancyhead[LE,RO]{\small\thepage}
\fancyhead[LO]{\small\textit{BioOS Causal Constitution}}
\fancyhead[RE]{\small\textit{Szőke L.-F.\ \textbullet\ Working Paper 2026}}
\renewcommand{\headrulewidth}{0.4pt}
\fancypagestyle{plain}{\fancyhf{}\fancyhead[LE,RO]{\small\thepage}\renewcommand{\headrulewidth}{0.4pt}}

% =============================================================
%  Section headings
% =============================================================
\titleformat{\section}
  {\large\bfseries\color{bioblue}}{\thesection.}{0.6em}{}
  [\vspace{-0.2em}\textcolor{bioblue}{\rule{\linewidth}{0.6pt}}]
\titleformat{\subsection}{\normalsize\bfseries\color{biogray}}{\thesubsection.}{0.5em}{}
\titleformat{\subsubsection}{\normalsize\itshape\color{biogray}}{\thesubsubsection.}{0.5em}{}

\setlength{\parskip}{3pt plus 1pt}
\setlength{\parindent}{1.2em}

% =============================================================
\begin{document}
% =============================================================

% ── Title page ───────────────────────────────────────────────
\begin{titlepage}
\begin{center}
  \vspace*{1.5cm}
  {\color{bioblue}\rule{\linewidth}{2pt}}\\[10pt]
  {\LARGE\bfseries\color{bioblue} BioOS Causal Constitution}\\[6pt]
  {\large\bfseries A Formally Verified Kernel-Level Causality Enforcement Primitive}\\[12pt]
  {\color{biogray}\rule{\linewidth}{0.5pt}}\\[8pt]
  {\large\color{biogray} BioOS Kauzális Alkotmány}\\[4pt]
  {\normalsize\color{biogray} Formálisan Verifikált Kernel-Szintű Kauzalitás-Érvényesítési Primitív}\\[14pt]
  {\color{bioblue}\rule{\linewidth}{0.5pt}}\\[22pt]
  {\large Szőke László-Ferenc}\\[4pt]
  {\normalsize Citrom Média LTD\,/\,LemonScript Laboratory}\\[2pt]
  {\small\texttt{hello@lemonscript.info}}\\[18pt]
  \begin{tabular}{rl}
    \textbf{Version:}   & Working Paper v0.1.0 \\
    \textbf{Date:}      & 2026-05-24 \\
    \textbf{Prototype:} & \bioURL \\
    \textbf{Z3 Proof:}  & \proofURL \\
    \textbf{Kernel:}    & Linux 6.1.0-48-cloud-amd64 · eBPF/LSM \\
    \textbf{Server:}    & GCP \texttt{asia-northeast1-b} (Tokyo) \\
  \end{tabular}\\[24pt]
  {\color{bioblue}\rule{\linewidth}{2pt}}\\[16pt]
  \begin{minipage}{0.9\linewidth}
    \small\itshape\color{biogray}
    Working paper submitted as a research artifact.
    All results are reproducible; the live prototype and Z3 proof pages
    are publicly accessible at the URLs above.
    Patent application: OSIM 20251221-2230.
  \end{minipage}
\end{center}
\end{titlepage}

% ── Abstracts ────────────────────────────────────────────────
\begin{abstract}
We present the \textbf{BioOS Causal Constitution}, a kernel-level security
primitive that enforces \emph{causal ancestry} rather than \emph{access permissions}.
The central claim is: a system call is permitted if and only if a verified causal
chain exists from a physical hardware interrupt (IRQ) to the call site, within a
bounded time window. In the absence of such a chain, protected resources do not
appear to be \emph{denied}---they appear to \emph{not exist}
(\texttt{ENODEV}, \texttt{ECONNREFUSED}).
We define the \texttt{.bio} specification language with a small-step operational
semantics, state the compiler soundness theorem for the eBPF/LSM backend, and
prove a bisimulation result between a JavaScript browser emulator
(\texttt{bio\_kernel\_emu}) and the kernel implementation (\textit{DCC Ring~0}).
All six security invariants are formally verified using the Z3 SMT solver
(violation $\Rightarrow$ \textsc{unsat}).
The implementation runs live on Linux~6.1 with six eBPF programs attached to
\texttt{raw\_tracepoint} and LSM hooks.
\end{abstract}

\vspace{6pt}
\begin{magyar}
\begin{abstract}
Bemutatjuk a \textbf{BioOS Kauzális Alkotmányt}, egy kernel-szintű biztonsági
primitívet, amely \emph{kauzális leszármazást} érvényesít a hagyományos
\emph{hozzáférési engedélyek} helyett. A fő állítás: egy rendszerhívás pontosan
akkor engedélyezett, ha egy ellenőrzött kauzális lánc létezik egy fizikai
hardvermegszakítástól (IRQ) a híváshelyig, egy korlátolt időablakban. Ilyen
lánc hiányában a védett erőforrások nem tiltottnak, hanem \emph{nem létezőnek}
tűnnek (\texttt{ENODEV}, \texttt{ECONNREFUSED}). Meghatározzuk a \texttt{.bio}
specifikációs nyelvet kis-lépéses operatív szemantikával, kimondják az eBPF/LSM
fordítói korrektségi tételt, és biszimulációs eredményt bizonyítunk egy JavaScript
böngészőemulátor (\texttt{bio\_kernel\_emu}) és a kernelimplementáció
(\textit{DCC Ring~0}) között. Mind a hat biztonsági invariánst formálisan
verifikáljuk a Z3 SMT-megoldóval (megsértés $\Rightarrow$ \textsc{unsat}).
Az implementáció élesben fut Linux~6.1 kernelen, hat eBPF-programmal csatolva
\texttt{raw\_tracepoint} és LSM hookokhoz.
\end{abstract}
\end{magyar}

\newpage
\tableofcontents
\newpage

% =============================================================
%  PART I — ENGLISH
% =============================================================
\part{English}

\section{The \texttt{.bio} Language: Syntax and Operational Semantics}

\subsection{Motivation}

The \texttt{.bio} language is a \textbf{Turing-incomplete} domain-specific language
for specifying safety-critical state machines. Turing-incompleteness is a design
requirement, not a limitation: it guarantees that the full state space is finite and
enumerable, making complete formal verification tractable.

\begin{mdframed}[style=notebox]
\textbf{Key property:} Every \texttt{.bio} program terminates. There are no
unbounded loops, no general recursion, no dynamic memory allocation. The compiler
can prove all reachable states at synthesis time.
\end{mdframed}

\subsection{Abstract Syntax (BNF)}

\vspace{4pt}
\begin{align*}
\nt{Program}   &\bnfdef \tm{CELL}\ \nt{Ident}\ \tm{\{}\ \nt{Interface}\ \nt{Invariants}\ \nt{States}\ \tm{\}} \\[3pt]
\nt{Interface}  &\bnfdef \tm{INTERFACE}\ \tm{\{}\ (\nt{Direction}\ \nt{Ident}\ \tm{:}\ \nt{Type}\ \tm{;})^*\ \tm{\}} \\
\nt{Direction}  &\bnfdef \tm{INPUT} \bnfor \tm{OUTPUT} \\
\nt{Type}       &\bnfdef \tm{BINARY} \bnfor \tm{INTEGER} \bnfor \tm{FLOAT}
                          \bnfor \tm{TIMESTAMP} \bnfor \tm{VECTOR2D}
                          \bnfor \tm{VECTOR3D} \bnfor \tm{ENUM}\ \tm{(}\ \nt{Ident}^+\ \tm{)} \\[3pt]
\nt{Invariants} &\bnfdef \tm{INVARIANTS}\ \tm{\{}\ \nt{Rule}^*\ \tm{\}} \\
\nt{Rule}       &\bnfdef \tm{RULE}\ \nt{Ident}\ \tm{:}\ \nt{Expr}\ \tm{;} \\[3pt]
\nt{States}     &\bnfdef \tm{STATES}\ \tm{\{}\ \nt{State}^*\ \tm{\}} \\
\nt{State}      &\bnfdef \tm{STATE}\ \nt{Ident}\ \nt{StateType}^?\ \tm{\{}\ \nt{Assign}^*\ \nt{Trans}^*\ \tm{\}} \\
\nt{StateType}  &\bnfdef \tm{TYPE}\ \tm{CRITICAL\_SHIELD} \\
\nt{Trans}      &\bnfdef \tm{TRANSITION}\ \tm{TO}\ \nt{Ident}\ \tm{IF}\ \nt{Expr}\ \tm{;} \\[3pt]
\nt{Expr}       &\bnfdef \nt{Atom} \bnfor \nt{Expr}\ \nt{BinOp}\ \nt{Expr}
                          \bnfor \nt{UnOp}\ \nt{Expr}
                          \bnfor \nt{Builtin}\ \tm{(}\ \nt{Expr}^+\ \tm{)} \\
\nt{BinOp}      &\bnfdef \tm{AND} \bnfor \tm{OR} \bnfor \tm{IMPLIES}
                          \bnfor \tm{==} \bnfor \tm{!=} \bnfor \tm{<}
                          \bnfor \tm{>} \bnfor \tm{+} \bnfor \tm{-} \\
\nt{UnOp}       &\bnfdef \tm{NOT} \\
\nt{Builtin}    &\bnfdef \tm{DISTANCE} \bnfor \tm{MAX\_DIFF}
                          \bnfor \tm{RATE\_OF\_CHANGE} \bnfor \tm{LIFESPAN} \\
\nt{Atom}       &\bnfdef \nt{Ident} \bnfor \nt{IntLit} \bnfor \nt{FloatLit}
                          \bnfor \tm{TRUE} \bnfor \tm{FALSE}
\end{align*}

\subsection{Small-Step Operational Semantics}

We define semantics over a \textbf{configuration} $\langle S,\,\sigma,\,\tau\rangle$:
$S \in \textit{States}$ (current state node),
$\sigma : \textit{Var} \to \textit{Val}$ (variable valuation),
$\tau \in \mathbb{N}$ (discrete time step).

\subsubsection{Expression evaluation}
\begin{align*}
\sem{n}_\sigma &= n \qquad \sem{x}_\sigma = \sigma(x) \\
\sem{e_1 \mathbin{\tm{AND}} e_2}_\sigma &= \sem{e_1}_\sigma \wedge \sem{e_2}_\sigma \qquad
\sem{e_1 \mathbin{\tm{IMPLIES}} e_2}_\sigma = \neg\sem{e_1}_\sigma \vee \sem{e_2}_\sigma \\
\sem{\tm{RATE\_OF\_CHANGE}(x)}_\sigma &= {|\sigma(x) - \sigma_{\mathrm{prev}}(x)|}/{\Delta\tau}
\end{align*}

\subsubsection{Invariant satisfaction}
\[
  \sigma \models P \;\iff\; \forall\, r_i \in P.\textit{Invariants}:\; \sem{r_i.\textit{expr}}_\sigma = \tm{TRUE}
\]

\subsubsection{State transition rule (TRANS)}
\[
\irule{
  \sem{\mathit{guard}}_\sigma = \tm{TRUE} \qquad
  \sigma' = \sigma[\text{assigns of }S'] \qquad
  \sigma' \models P
}{
  \langle S,\,\sigma,\,\tau\rangle \;\longrightarrow\; \langle S',\,\sigma',\,\tau{+}1\rangle
  \quad \text{where } (\tm{TRANSITION TO}\;S'\;\tm{IF}\;\mathit{guard})\in S
}{Trans}
\]
If no guard is satisfied the system self-loops in $S$.
A \tm{CRITICAL\_SHIELD} state is a \textbf{deadlock sink} by design.

\subsubsection{Invariant violation (APOPTOSIS)}
\[
\irule{
  \sigma' \not\models P
}{
  \langle S,\,\sigma,\,\tau\rangle \;\longrightarrow\;
  \langle S_0,\,\sigma_{\mathrm{snap}},\,\tau{+}1\rangle \quad
  S_0 = \text{initial state},\; \sigma_{\mathrm{snap}} = \text{last valid snapshot}
}{Apoptosis}
\]
Any invariant violation atomically reverts to the last valid snapshot,
mirroring \texttt{bpf\_map\_update\_elem(BPF\_EXIST)} atomicity.

\subsection{Turing-Incompleteness Theorem}

\begin{theorem}[Termination]\label{thm:termination}
Every \texttt{.bio} program halts.
\end{theorem}
\begin{proof}[Proof sketch]
The state space is finite (declared at compile time). The transition relation is
deterministic (at most one guard true at any time). No cycles through
\tm{CRITICAL\_SHIELD} states exist. Z3 enumerates all reachable states at
synthesis time. \qed
\end{proof}

\begin{corollary}
The full state space is bounded by $|\textit{States}|\times|\textit{Val}^{\textit{Var}}|$,
finite for bounded integer/boolean types. Complete invariant verification is decidable.
\end{corollary}

% ─────────────────────────────────────────────────────────────
\section{BioOS Ring~0 Guard as a \texttt{.bio} Program}

\subsection{Full Specification}

\begin{lstlisting}[style=biostyle,caption={DCCRing0Guard -- canonical \texttt{.bio} specification}]
CELL DCCRing0Guard {

  INTERFACE {
    INPUT  irq_event:    BINARY;   // Hardware IRQ? (input_event, EV_KEY)
    INPUT  op_class:     INTEGER;  // Token bitmask: OP_WRITE=1 OP_NET=2 OP_EXEC=4
    INPUT  file_axiom:   INTEGER;  // file_axiom_map (0=BLOCK, 255=ANY)
    INPUT  token_age_ms: FLOAT;    // Token age in milliseconds
    INPUT  tx_state:     INTEGER;  // TX state (0=IDLE 1=LOCKED 2=STAGED)
    INPUT  same_inode:   BINARY;   // Same inode as bound tx?
    OUTPUT verdict:      INTEGER;  // 1=SAT 2=NO_TOKEN 3=EXPIRED 11=ROLLBACK
  }

  INVARIANTS {
    // I1: No autonomous write -- the central causal requirement
    RULE no_autonomous_write:
        (irq_event == FALSE) IMPLIES (verdict != 1);

    // I2: Hard 500 ms causality window
    RULE causality_window: token_age_ms < 500.0;

    // I3: Protected files immutable regardless of token state
    RULE blocked_file_immutable:
        (file_axiom == 0) IMPLIES (verdict != 1);

    // I4: Op_class mismatch -- OP_WRITE cannot grant OP_NET
    RULE op_class_match_required:
        ((file_axiom != 255) AND ((file_axiom AND op_class) == 0))
            IMPLIES (verdict != 1);

    // I5: TOCTOU -- staged TX + different inode = rollback, always
    RULE toctou_rollback:
        ((tx_state == 2) AND (same_inode == FALSE)) IMPLIES (verdict == 11);

    // I6: SAT requires complete causal chain
    RULE sat_requires_causal_chain:
        (verdict == 1) IMPLIES (irq_event == TRUE AND token_age_ms < 500.0);
  }

  STATES {
    STATE TX_IDLE {
      verdict = 2;  // NO_TOKEN
      TRANSITION TO TX_LOCKED IF irq_event == TRUE AND token_age_ms < 500.0;
    }
    STATE TX_LOCKED {
      TRANSITION TO TX_STAGED
          IF file_axiom != 0
         AND (file_axiom == 255 OR (file_axiom AND op_class) != 0);
      TRANSITION TO TX_IDLE IF token_age_ms >= 500.0;
    }
    STATE TX_STAGED {
      verdict = 1;  // SAT
      TRANSITION TO TX_STAGED IF same_inode == TRUE AND token_age_ms < 500.0;
      TRANSITION TO TX_IDLE   IF same_inode == FALSE;
      TRANSITION TO TX_IDLE   IF token_age_ms >= 500.0;
    }
    STATE TX_COMMITTED TYPE CRITICAL_SHIELD { verdict = 1; }
  }
}
\end{lstlisting}

\subsection{State Machine Diagram}

\begin{figure}[H]
\centering
\begin{tikzpicture}[
  node distance=3.4cm,
  st/.style={rectangle,rounded corners=6pt,draw=bioblue,line width=1.2pt,
             fill=lightbluebg,text width=2.7cm,align=center,
             font=\small\ttfamily,inner sep=7pt},
  sink/.style={rectangle,rounded corners=6pt,draw=biored,line width=1.5pt,
               fill=red!8,text width=2.7cm,align=center,
               font=\small\ttfamily,inner sep=7pt},
  >=Stealth, auto, font=\footnotesize,
]
\node[st]   (idle)   {TX\_IDLE\\[2pt]{\color{biogray}verdict=2}};
\node[st]   (locked) [right=of idle]   {TX\_LOCKED\\[2pt]{\color{biogray}awaiting write}};
\node[st]   (staged) [right=of locked] {TX\_STAGED\\[2pt]{\color{biogreen}verdict=1}};
\node[sink] (comm)   [below=2cm of staged] {TX\_COMMITTED\\[2pt]{\color{biored}CRITICAL\_SHIELD}};

\draw[->] (idle)   -- node[above,align=center]{irq=TRUE\\age$<$500ms} (locked);
\draw[->] (locked) -- node[above,align=center]{axiom OK\\op match}    (staged);
\draw[->] (staged) -- node[right]{commit} (comm);
\draw[->,bend right=30] (locked) to node[below,align=center]{timeout} (idle);
\draw[->,bend right=32] (staged) to node[below,align=center]{TOCTOU\\timeout} (idle);
\draw[->] (staged) to[out=70,in=110,loop] node[above]{same\_inode} (staged);
\end{tikzpicture}
\caption{DCCRing0Guard state machine. \texttt{TX\_COMMITTED} is a deadlock sink.}
\end{figure}

% ─────────────────────────────────────────────────────────────
\section{Threat Model}

\subsection{Attacker Capabilities}

$\mathcal{A}$ is a Ring~3 adversary. \textbf{$\mathcal{A}$ can:}
arbitrary Ring~3 code execution (shellcode, ROP, thread hijack);
fork, thread, all POSIX syscalls; attempt file writes, network connects, exec;
\texttt{prctl(PR\_SET\_NAME)} spoofing; headless operation; software timers/cron.
\textbf{$\mathcal{A}$ cannot:}
generate hardware \texttt{EV\_KEY} events without physical input;
set \texttt{isTrusted=true} on synthetic browser events;
modify loaded eBPF code;
extend a \texttt{CausalToken} past \texttt{CAUSALITY\_WINDOW\_NS}.

\subsection{Trust Boundary Map}

\begin{figure}[H]
\centering
\begin{tikzpicture}[font=\small]
  \node[draw=biored,thick,rounded corners,fill=red!5,
        text width=9cm,align=left,inner sep=10pt] (outer) {
    \textbf{\color{biored}Outside model ($\mathcal{A}$ controls)}\\[5pt]
    \quad Ring 3 userspace: processes, threads, memory\\
    \quad Network stack (outbound connection attempts)\\
    \quad Software timers, cron, scripted automation
  };
  \node[draw=bioblue,thick,rounded corners,fill=lightbluebg,
        text width=7.2cm,align=left,inner sep=10pt,
        below=0.4cm of outer] (inner) {
    \textbf{\color{bioblue}Inside model (DCC Ring~0)}\\[5pt]
    \quad \texttt{global\_state\_map} (BPF\_MAP\_TYPE\_ARRAY)\\
    \quad \texttt{token\_map} (BPF\_MAP\_TYPE\_HASH, per-PID)\\
    \quad \texttt{tx\_map} (transaction state machine)\\
    \quad \texttt{file\_axiom\_map} (protected file set)\\
    \quad IRQ source: \texttt{raw\_tp/input\_event}
  };
  \draw[->,thick,biogray] (outer.south) --
    node[right]{\textit{LSM hooks (kernel boundary)}} (inner.north);
\end{tikzpicture}
\caption{Trust boundary. $\mathcal{A}$ controls the outer zone.}
\end{figure}

\subsection{Attack Scenarios and Formal Refutations}

\begin{longtable}{@{}p{3.2cm}p{3.8cm}p{5.4cm}@{}}
\toprule
\textbf{Attack} & \textbf{Mechanism} & \textbf{Formal refutation} \\ \midrule
\endhead \bottomrule \endfoot
Autonomous write (ransomware)
  & File write without user interaction
  & I1: $\mathit{irq}{=}F \Rightarrow \mathit{verdict}{\neq}1$ (UNSAT, Z3) \\[3pt]
TOCTOU pivot
  & Write A, redirect to B (different inode)
  & I5: $\mathit{tx}{=}2 \wedge \neg\mathit{same\_inode} \Rightarrow \mathit{verdict}{=}11$ \\[3pt]
Op\_class confusion (camera${\to}$net)
  & OP\_WRITE token for network socket
  & I4: $(\mathit{axiom}{\neq}255) \wedge (\mathit{axiom}\,\&\,\mathit{op}){=}0 \Rightarrow \mathit{verdict}{\neq}1$ \\[3pt]
Token replay
  & Expired or consumed token reuse
  & I2: $\mathit{age}{\geq}500 \Rightarrow \mathit{verdict}{=}3$; per-use consumed flag \\[3pt]
Blocked file bypass
  & Write to \texttt{--protect}-ed file
  & I3: $\mathit{axiom}{=}0 \Rightarrow \mathit{verdict}{\neq}1$ \\[3pt]
Headless server
  & No \texttt{/dev/input} device
  & \texttt{raw\_tp/input\_event} never fires $\Rightarrow$ \texttt{INERT} invariant \\[3pt]
Thread hijack
  & Inject into process with valid token
  & Token per-PID in \texttt{token\_map}; fork inheritance max 3 generations \\
\end{longtable}

% ─────────────────────────────────────────────────────────────
\section{Compiler Backend and Isomorphism Theorem}

\subsection{Translation Overview}

\begin{table}[H]
\centering
\caption{Compilation mapping: \texttt{.bio} $\to$ eBPF/LSM}
\begin{tabular}{@{}ll@{}}
\toprule
\textbf{\texttt{.bio} construct} & \textbf{eBPF/LSM construct} \\ \midrule
\texttt{CELL} & BPF program set, single object \\
\texttt{INPUT x: BINARY} & BPF map entry or \texttt{ctx} field \\
\texttt{OUTPUT verdict: INTEGER} & BPF program \texttt{return} value \\
\texttt{RULE r: e} & \texttt{if (!eval(e,ctx)) return -ENODEV;} \\
\texttt{STATE S \{...\}} & \texttt{global\_state\_map} + transitions \\
\texttt{TRANSITION TO S' IF g} & \texttt{bpf\_map\_update\_elem(\&state, ...)} \\
\texttt{TYPE CRITICAL\_SHIELD} & no outgoing transitions compiled \\
Apoptosis & \texttt{bpf\_map\_update\_elem(STATE\_INERT)} \\
\bottomrule
\end{tabular}
\end{table}

\subsection{Compiler Soundness}

\begin{theorem}[Compiler Soundness]\label{thm:soundness}
Let $P$ be a valid \texttt{.bio} program, $\mathcal{C}(P)$ its compiled eBPF/LSM
implementation. Then:
\[
  \forall\,\sigma:\quad \sigma \in \sem{\mathcal{C}(P)} \;\Longrightarrow\; \sigma \in \sem{P}
\]
The compiled kernel code never permits a trace the \texttt{.bio} specification forbids.
\end{theorem}
\begin{proof}[Proof sketch]
(1)~Invariant translation is conservative: each \texttt{RULE} compiles to an
\texttt{if}-check \emph{before} any state transition.
(2)~State transitions are atomic via \texttt{bpf\_map\_update\_elem}.
(3)~\texttt{CRITICAL\_SHIELD} states emit no transition code.
(4)~Token age is kernel-clock enforced via monotonic \texttt{bpf\_ktime\_get\_ns()}.
Full mechanization in Isabelle/HOL is future work. \qed
\end{proof}

\subsection{Bisimulation}

\begin{definition}[LTS]\label{def:lts}
$\mathrm{LTS}(X) = (Q_X, \Sigma, {\to_X})$ where $Q_X$ are configurations
$\langle S,\sigma,\tau\rangle$, $\Sigma$ is the event alphabet (IRQ, resource
request, verdict), ${\to_X}$ the transition relation.
\end{definition}

\begin{definition}[Bisimulation]\label{def:bisim}
$R \subseteq Q_\mathrm{JS} \times Q_K$ is a \emph{bisimulation} if:
\[
\forall\,(q_\mathrm{JS},q_K)\in R:\begin{cases}
  q_\mathrm{JS}\xrightarrow{a}q'_\mathrm{JS} \Rightarrow \exists q'_K: q_K\xrightarrow{a}q'_K \wedge (q'_\mathrm{JS},q'_K)\in R \\
  q_K\xrightarrow{a}q'_K \Rightarrow \exists q'_\mathrm{JS}: q_\mathrm{JS}\xrightarrow{a}q'_\mathrm{JS} \wedge (q'_\mathrm{JS},q'_K)\in R
\end{cases}
\]
\end{definition}

\begin{theorem}[Bisimulation]\label{thm:bisim}
$\mathrm{LTS}(\texttt{bio\_kernel\_emu})$ and $\mathrm{LTS}(\textit{DCC Ring~0})$
are bisimilar under the correspondence in Table~\ref{tab:bisim}.
\end{theorem}

\begin{table}[H]
\centering
\caption{Bisimulation correspondence\label{tab:bisim}}
\begin{tabular}{@{}ll@{}}
\toprule
\textbf{bio\_kernel\_emu (JS)} & \textbf{DCC Ring~0 (eBPF)} \\ \midrule
\texttt{STATE\_INERT}              & \texttt{global\_state\_map[0] = 0} \\
\texttt{STATE\_ACTIVE}             & \texttt{global\_state\_map[0] = 1} \\
\texttt{mousedown.isTrusted=true}  & \texttt{raw\_tp/input\_event, EV\_KEY, value=1} \\
\texttt{CausalToken\{age<200ms\}}  & \texttt{causal\_token\{age<CAUSALITY\_WINDOW\_NS\}} \\
\texttt{Z3Gatekeeper.verify()=SAT} & \texttt{dcc\_axiom\_validator() return 0} \\
\texttt{Z3Gatekeeper.verify()=UNSAT}& \texttt{dcc\_axiom\_validator() return -ENODEV} \\
\texttt{stateManager.rollback()}   & \texttt{bpf\_map\_update\_elem(STATE\_INERT)} \\
\bottomrule
\end{tabular}
\end{table}

\begin{proof}[Proof sketch]
Both systems derive from the same \texttt{.bio} source. The Z3 verification
(\texttt{verify\_isomorphism.py}) proves all 6 invariants hold in both systems
simultaneously. No trace violates an invariant in one without violating it in the
other (safety bisimulation). Liveness follows from the shared state machine. \qed
\end{proof}

\begin{remark}
The causality window (200\,ms JS vs.\ 500\,ms kernel) is an implementation
parameter, not part of the invariant structure.
\end{remark}

% ─────────────────────────────────────────────────────────────
\section{Experimental Results}

\subsection{Setup}

\begin{table}[H]
\centering
\begin{tabular}{@{}ll@{}}
\toprule
\textbf{Component} & \textbf{Value} \\ \midrule
Kernel             & Linux 6.1.0-48-cloud-amd64 \\
Machine            & GCP \texttt{asia-northeast1-b} (Tokyo), 2\,vCPU 4\,GB \\
eBPF toolchain     & libbpf, CO-RE, clang\,14 \\
Z3 (server)        & 4.16.0 (python3-z3) \\
Z3 (local)         & 4.15.4 \\
Browser emulator   & bio\_kernel\_emu v0.8.1, V8 \\
\bottomrule
\end{tabular}
\end{table}

\subsection{Z3 Isomorphism Verification --- 6/6 PASS}

\begin{lstlisting}[style=shellstyle,caption={Z3 verifier output (Tokyo server, Z3 4.16.0)}]
=== BioOS Causal Constitution -- Z3 Isomorphism Verifier ===
    bio_kernel_emu (JS)  <->  DCC Ring 0 (eBPF/LSM)
    Tokyo server  kernel 6.1.0-48-cloud-amd64

-- Group I: Physical Causality --------------------------------

  [PASS] I1 -- STATE_ACTIVE requires isTrusted=true IRQ
         -> BPF: dcc_causality_monitor  [raw_tp/input_event]
         -> Headless server: no /dev/input -> INERT invariant

  [PASS] I2 -- Token validity window (CAUSALITY_WINDOW_MS = 200 ms)
         -> BPF: dcc_token_validator  [token_map TTL]

-- Group II: Error Semantics ----------------------------------

  [PASS] I3 -- INERT file access returns ENODEV (errno 19)
         -> BPF: dcc_axiom_validator  [lsm/file_permission]

  [PASS] I4 -- INERT network access returns ECONNREFUSED (errno 111)
         -> BPF: dcc_network_guard  [lsm/socket_connect]

-- Group III: Integrity & Scope -------------------------------

  [PASS] I5 -- TOCTOU: staged TX + inode mismatch -> result = -1
         -> BPF: dcc_toctou_guard  [lsm/file_permission, inode]

  [PASS] I6 -- Op_class scope: granted iff token_op & req_op != 0
         -> BPF: dcc_scope_checker  [all LSM hooks]

=== Result: 6/6 invariants verified  ISOMORPHISM PROVEN ===
\end{lstlisting}

\noindent Z3 Python encoding (I1 --- violation $\Rightarrow$ UNSAT):

\begin{lstlisting}[style=pythonstyle]
from z3 import *
STATE_ACTIVE = 1
s = Solver()
state, is_trusted = Int('state'), Bool('is_trusted')
s.add(Implies(state == STATE_ACTIVE, is_trusted == True))   # axiom
s.add(And(state == STATE_ACTIVE, is_trusted == False))       # violation
assert s.check() == unsat   # -> PASS
\end{lstlisting}

\subsection{Live Kernel --- 6 eBPF Programs}

\begin{lstlisting}[style=shellstyle,caption={Live eBPF programs on Tokyo server}]
$ sudo bpftool prog list
323: raw_tracepoint  name dcc_causality_monitor  tag 7061cd22c7437b9c  gpl
324: raw_tracepoint  name dcc_fork_inherit       tag 55cfd16ceb2b3666  gpl
325: lsm             name dcc_axiom_validator    tag ca6a3e97b0d9cda2  gpl
326: lsm             name dcc_read_guard         tag edf568e254142683  gpl
327: lsm             name dcc_network_guard      tag 0d93cfdec6d3c3f7  gpl
328: lsm             name dcc_exec_guard         tag 11b02161c7c6b71a  gpl
\end{lstlisting}

\subsection{Attack Scenario Tests}

\begin{table}[H]
\centering
\caption{Test results (\texttt{test\_phase89.sh}, Tokyo 2026-05-21)}
\begin{tabular}{@{}lllc@{}}
\toprule
\textbf{Test} & \textbf{Scenario} & \textbf{Expected} & \textbf{Result} \\ \midrule
T1 & Autonomous write (headless, no IRQ) & NO\_TOKEN        & \checkmark \\
T2 & TOCTOU: write A then B same token   & TX\_ROLLBACK     & \checkmark \\
T3 & Legitimate multi-write same file    & SAT              & \checkmark \\
T4 & Write \texttt{--protect}-ed file (\texttt{axiom=0}) & AXIOM\_MISMATCH & \checkmark \\
T5 & Protected file with valid token     & AXIOM\_MISMATCH  & \checkmark \\
T6 & Phase 7 regression (\texttt{prove\_ring0.sh}) & 11/11 PASS & \checkmark \\
\midrule
\multicolumn{3}{l}{\textbf{Total} (\texttt{run\_proofs.sh})} & \textbf{7/7 PASS} \\
\bottomrule
\end{tabular}
\end{table}

% ─────────────────────────────────────────────────────────────
\section{Discussion}

\subsection{Relation to Existing Work}

\textbf{Causal closure / information flow.}
DCC enforces \emph{physical} causality (hardware IRQ as root) rather than
information-theoretic causality (Clarkson \& Schneider, JCS 2010).

\textbf{eBPF security.}
Prior work (KSplit, BPF verifier) targets individual program correctness.
We use BPF as a \emph{carrier} for policy derived from a Turing-incomplete DSL.

\textbf{MAC/DAC vs.\ causal enforcement.}
MAC answers ``is this principal allowed?''; DCC answers ``does this operation
have a physical causal ancestor?'' The models are orthogonal and composable.

\subsection{Limitations and Future Work}

\begin{enumerate}[topsep=2pt,itemsep=2pt]
  \item Compiler soundness mechanization (Isabelle/HOL or Lean~4).
  \item Bisimulation liveness direction (BPF map update timing analysis).
  \item Comm-whitelist spoofing via \texttt{prctl}: switch to cgroup-based whitelisting.
  \item Systemd service for DCC Ring~0 auto-load (Phase~11).
  \item Formal \texttt{.bio} encoding in TLA+ or Coq.
\end{enumerate}

% ─────────────────────────────────────────────────────────────
\section{Conclusion}

The BioOS Causal Constitution contributes:
(1) \texttt{.bio} small-step semantics with Turing-incompleteness guarantee;
(2) DCCRing0Guard: six security invariants in a concrete \texttt{.bio} program;
(3) explicit threat model with six attack scenarios formally refuted via Z3;
(4) compiler soundness theorem and bisimulation (JS emulator $\leftrightarrow$ kernel);
(5) live experimental validation: 6/6 Z3 PASS, 7/7 system tests on Linux~6.1.

The prototype is publicly accessible and all proofs are reproducible.

\appendix

\section{Reproducibility}

\begin{lstlisting}[style=shellstyle]
# Local Z3 verification (requires: pip install z3-solver)
python bio_kernel_emu/tests/verify_isomorphism.py
# Expected: 6/6 PASS

# On Tokyo server
ssh lszok@34.146.249.102 \
  "python3 /home/lszok/bio_kernel_emu_tests/verify_isomorphism.py"
# Expected: 6/6 PASS

# Kernel tests
ssh lszok@34.146.249.102 \
  "cd /home/lszok/dcc_ebpf/src && sudo bash tests/run_proofs.sh"
# Expected: 7/7 PASS
\end{lstlisting}

Interactive demo: \bioURL\quad Z3 proof page: \proofURL

\section{Artifact Locations}

\begin{tabular}{@{}ll@{}}
\toprule
\textbf{Artifact} & \textbf{Location} \\ \midrule
\texttt{.bio} specification   & \texttt{DCC\_RING0/spec/dcc\_ring0.bio} \\
eBPF source                   & \texttt{DCC\_RING0/src/dcc\_core.bpf.c} \\
Z3 isomorphism verifier       & \texttt{bio\_kernel\_emu/tests/verify\_isomorphism.py} \\
Kernel test suite             & \texttt{DCC\_RING0/tests/run\_proofs.sh} \\
JS emulator                   & \texttt{bio\_kernel\_emu/demo/} \\
Causal constitution spec      & \texttt{bio\_kernel\_emu/spec/causal\_constitution.md} \\
\bottomrule
\end{tabular}

\newpage

% =============================================================
%  PART II — MAGYAR VÁLTOZAT
% =============================================================
\begin{magyar}
\part{Magyar változat}

\section{A \texttt{.bio} Nyelv: Szintaxis és Operatív Szemantika}

\subsection{Motiváció}

A \texttt{.bio} nyelv egy \textbf{Turing-hiányos} tartomány-specifikus nyelv
biztonsági szempontból kritikus állapotgépek specifikálásához.
A Turing-hiányosság tervezési követelmény, nem korlát: garantálja, hogy a
teljes állapottér véges és felsorolható, ami teljes formális verifikációt tesz lehetővé.

\begin{mdframed}[style=notebox]
\textbf{Kulcstulajdonság:} Minden \texttt{.bio} program megáll.
Nincsenek korlátlan ciklusok, általános rekurzió, dinamikus memóriaallokáció.
A fordító szintézis idején be tudja bizonyítani az összes elérhető állapotot.
\end{mdframed}

\subsection{Absztrakt Szintaxis (BNF)}

A grammatika megegyezik az angol változatéval (ld.\ 1.2~szakasz). A nemterminálisok
dőlt betűvel, a terminálisok írógép-betűtípussal szerepelnek.

\subsection{Kis-lépéses Operatív Szemantika}

A szemantikát egy \textbf{konfigurációra} $\langle S,\,\sigma,\,\tau\rangle$
definiáljuk:
$S \in \textit{Állapotok}$ (aktuális állapotcsúcs),
$\sigma : \textit{Változó} \to \textit{Érték}$ (változóértékelés),
$\tau \in \mathbb{N}$ (diszkrét időlépés, monoton növekvő).

\subsubsection{Invariáns-teljesítés}
\[
  \sigma \models P \;\iff\; \forall\, r_i \in P.\textit{Invariánsok}:\;
  \sem{r_i.\textit{kif}}_\sigma = \tm{TRUE}
\]

\subsubsection{Állapotátmeneti szabály}
\[
\irule{
  \sem{\mathit{feltétel}}_\sigma = \tm{TRUE} \qquad
  \sigma' = \sigma[\text{S' értékadásai}] \qquad
  \sigma' \models P
}{
  \langle S,\,\sigma,\,\tau\rangle \;\longrightarrow\;
  \langle S',\,\sigma',\,\tau{+}1\rangle
}{Trans}
\]

Ha egyetlen feltétel sem teljesül, a rendszer önhurokba esik. A
\tm{CRITICAL\_SHIELD} típusú állapot holtponti elnyelő: nincsenek kimenő átmenetek.

\subsubsection{Invariánssértés — Apoptózis-szabály}
\[
\irule{
  \sigma' \not\models P
}{
  \langle S,\,\sigma,\,\tau\rangle \;\longrightarrow\;
  \langle S_0,\,\sigma_{\mathrm{pillanat}},\,\tau{+}1\rangle
}{Apoptózis}
\]

Nincsen részleges hiba: bármely invariánssértés atomikusan visszaállítja az
utolsó érvényes pillanatképet.

\subsection{Turing-leállási Tétel}

\begin{tetel}[Leállás]\label{thm:leallas}
Minden \texttt{.bio} program megáll.
\end{tetel}
\begin{proof}[Bizonyítás vázlata]
Az állapottér véges. Az átmeneti reláció determinisztikus.
Nincsenek \tm{CRITICAL\_SHIELD} állapoton átmenő ciklusok.
A Z3-verifikátor szintézis idején felsorolja az összes elérhető állapotot. \qed
\end{proof}

\begin{kovetkezmeny}
A teljes állapottér $|\textit{Állapotok}|\times|\textit{Érték}^{\textit{Változó}}|$
értékkel korlátozott; a teljes invariáns-verifikáció eldönthető.
\end{kovetkezmeny}

% ─────────────────────────────────────────────────────────────
\section{A BioOS Ring~0 Őr mint \texttt{.bio} Program}

\subsection{Teljes Specifikáció}

Az annotált \texttt{.bio} specifikáció az angol változat 2.1~szakaszában
található (1.~lista). A hat invariáns biztonsági szemantikáját az alábbi
táblázat foglalja össze:

\begin{table}[H]
\centering
\caption{A hat kauzális invariáns}
\begin{tabular}{@{}llp{6cm}@{}}
\toprule
\textbf{Inv.} & \textbf{Azonosító} & \textbf{Tartalom} \\ \midrule
I1 & \texttt{no\_autonomous\_write}   & IRQ nélkül nincs SAT \\
I2 & \texttt{causality\_window}       & token\_age < 500\,ms \\
I3 & \texttt{blocked\_file\_immutable}& file\_axiom=0 esetén nincs SAT \\
I4 & \texttt{op\_class\_match}        & bitmask-egyezés kötelező \\
I5 & \texttt{toctou\_rollback}        & más inode = visszagörgetés \\
I6 & \texttt{sat\_requires\_chain}    & SAT $\Rightarrow$ teljes kauzális lánc \\
\bottomrule
\end{tabular}
\end{table}

\subsection{Állapotgép}

Az állapotgép négy csúcsból áll (ld.\ az angol változat 2.~szakaszának ábráját):
\texttt{TX\_IDLE} (nincs token),
\texttt{TX\_LOCKED} (token érvényes, első írásra vár),
\texttt{TX\_STAGED} (inode kötött, SAT érvényes),
\texttt{TX\_COMMITTED} (\tm{CRITICAL\_SHIELD}, terminális).

% ─────────────────────────────────────────────────────────────
\section{Fenyegetési Modell}

\subsection{A Támadó Képességei}

A $\mathcal{A}$ támadó Ring~3-ban tetszőleges kódot futtathat, folyamatokat
forkolhat, minden POSIX-rendszerhívást használhat, folyamatnevet hamisíthat,
és szoftveres időzítőket ütemezhet. \textbf{Nem képes} fizikai hardver nélkül
\texttt{EV\_KEY} eseményt generálni, \texttt{isTrusted=true} értéket
programból beállítani, betöltött eBPF-kódot módosítani, vagy a
\texttt{CausalToken}-t meghosszabbítani.

\subsection{Támadási Forgatókönyvek}

\begin{longtable}{@{}p{3.2cm}p{3.8cm}p{5.4cm}@{}}
\toprule
\textbf{Támadás} & \textbf{Mechanizmus} & \textbf{Formális cáfolat} \\
\midrule \endhead \bottomrule \endfoot
Autonóm írás (zsarolóvírus)
  & Fájlírás felhasználói beavatkozás nélkül
  & I1: $\mathit{irq}{=}H \Rightarrow \mathit{verdict}{\neq}1$ (Z3 UNSAT) \\[3pt]
TOCTOU pivot
  & A fájl megnyitása, B fájlba írás
  & I5: $\mathit{tx}{=}2 \wedge \neg\mathit{same} \Rightarrow \mathit{verdict}{=}11$ \\[3pt]
Op\_class felcserélés
  & OP\_WRITE token hálózati sockethez
  & I4: bitmask-ütközés $\Rightarrow$ megtagadás \\[3pt]
Token visszajátszás
  & Lejárt token újrahasználata
  & I2: $\mathit{kor}{\geq}500\,\mathrm{ms} \Rightarrow$ EXPIRED \\[3pt]
Védett fájl megkerülése
  & \texttt{--protect}-elt fájlba írás
  & I3: axiom=0 $\Rightarrow$ megtagadás \\[3pt]
Fej nélküli szerver
  & Nincs \texttt{/dev/input}
  & Strukturálisan lehetetlen; \texttt{INERT} invariáns \\[3pt]
Szálrablás
  & Injektálás érvényes tokennel rendelkező folyamatba
  & PID-enkénti \texttt{token\_map}; max 3 generáció \\
\end{longtable}

% ─────────────────────────────────────────────────────────────
\section{A Fordítói Háttér és az Izomorfizmus-tétel}

\subsection{Fordítói Korrektség}

\begin{tetel}[Fordítói korrektség]\label{thm:korrektség}
Legyen $P$ érvényes \texttt{.bio} program, $\mathcal{C}(P)$ a lefordított
eBPF/LSM implementáció. Ekkor:
\[
  \forall\,\sigma:\quad \sigma \in \sem{\mathcal{C}(P)} \;\Longrightarrow\; \sigma \in \sem{P}
\]
\end{tetel}
\begin{proof}[Bizonyítás vázlata]
(1)~Az invariáns-fordítás konzervatív: minden \texttt{RULE} átmenet előtt
\texttt{if}-ellenőrzéssé fordul.
(2)~Az állapotátmenet atomi (\texttt{bpf\_map\_update\_elem}).
(3)~\tm{CRITICAL\_SHIELD} állapotokhoz nem fordulnak kimenő átmenetek.
(4)~A token érvényességét a kernelóra érvényesíti. \qed
\end{proof}

\subsection{Biszimiláció}

\begin{definicio}[Biszimiláció]\label{def:biszim_hu}
Az $R \subseteq Q_\mathrm{JS} \times Q_K$ reláció \emph{biszimiláció} a
\texttt{bio\_kernel\_emu} (JS) és a DCC Ring~0 (K) között, ha:
\[
\forall\,(q_\mathrm{JS},q_K)\in R:\begin{cases}
  q_\mathrm{JS}\xrightarrow{a}q'_\mathrm{JS} \Rightarrow \exists q'_K: q_K\xrightarrow{a}q'_K \wedge (q'_\mathrm{JS},q'_K)\in R \\
  q_K\xrightarrow{a}q'_K \Rightarrow \exists q'_\mathrm{JS}: q_\mathrm{JS}\xrightarrow{a}q'_\mathrm{JS} \wedge (q'_\mathrm{JS},q'_K)\in R
\end{cases}
\]
\end{definicio}

\begin{tetel}[Biszimiláció]\label{thm:biszim_hu}
$\mathrm{LTS}(\texttt{bio\_kernel\_emu})$ és $\mathrm{LTS}(\textit{DCC Ring~0})$
biszimilárisak az angol változat 4.~táblázatában megadott megfeleltetés alatt.
\end{tetel}

\begin{megjegyzes}
A kauzalitás-ablak (200\,ms JS, 500\,ms kernel) implementációs paraméter,
nem az invariánsstruktúra része.
\end{megjegyzes}

% ─────────────────────────────────────────────────────────────
\section{Kísérleti Eredmények}

\subsection{Z3 Izomorfizmus-ellenőrzés --- 6/6 PASS}

Minden $I_n$ invariánsnál a Z3 \texttt{Solver} a rendszeraxiómákkal inicializált,
majd az invariáns tagadása kerül hozzáadásra. A $\texttt{check()} = \textsc{unsat}$
eredmény bizonyítja, hogy a sértés strukturálisan lehetetlen. A teljes kimenet
és Python-kód az angol változat 5.~szakaszában található.

\subsection{Támadási Tesztek --- 7/7 PASS}

\begin{table}[H]
\centering
\caption{Teszteredmények (tokiói szerver, 2026-05-21)}
\begin{tabular}{@{}lllc@{}}
\toprule
\textbf{Teszt} & \textbf{Forgatókönyv} & \textbf{Várt} & \textbf{Eredmény} \\ \midrule
T1 & Autonóm írás (fej nélküli szerver)       & NO\_TOKEN       & \checkmark \\
T2 & TOCTOU: A fájl, majd B fájl              & TX\_ROLLBACK    & \checkmark \\
T3 & Törvényes többszöri írás                 & SAT             & \checkmark \\
T4 & Védett fájl (\texttt{axiom=0})           & AXIOM\_MISMATCH & \checkmark \\
T5 & Védett fájl érvényes tokennel            & AXIOM\_MISMATCH & \checkmark \\
T6 & 7.\ fázis regresszió                     & 11/11 PASS      & \checkmark \\
\midrule
\multicolumn{3}{l}{\textbf{Összesen}} & \textbf{7/7 PASS} \\
\bottomrule
\end{tabular}
\end{table}

% ─────────────────────────────────────────────────────────────
\section{Értékelés és Összefoglalás}

\subsection{Kapcsolat a Meglévő Munkákkal}

A DCC-modell \emph{fizikai} kauzalitást érvényesít (IRQ mint gyökér), nem
információelméleti kauzalitást. A BPF-et egy Turing-hiányos DSL-ből levezetett,
formálisan specifikált biztonsági politika hordozójaként használjuk. A DCC
ortogonális a hagyományos MAC/DAC modellekhez (SELinux, AppArmor) és
azokkal kompozálható.

\subsection{Korlátok és Jövőbeli Munka}

\begin{enumerate}[topsep=2pt,itemsep=2pt]
  \item Gépesített fordítói korrektség (Isabelle/HOL vagy Lean~4).
  \item Biszimiláció élőség iránya.
  \item Comm-névhamisítás elleni védelem: cgroup-alapú fehérlistázás.
  \item Systemd-szolgáltatás a DCC Ring~0 automatikus betöltésére (11.~fázis).
  \item Formális \texttt{.bio} szemantika TLA+-ban vagy Coq-ban.
\end{enumerate}

\subsection{Összefoglalás}

A BioOS Kauzális Alkotmány öt fő hozzájárulást tartalmaz:
(1) \texttt{.bio} kis-lépéses szemantika Turing-leállási garanciával;
(2) DCCRing0Guard: hat biztonsági invariáns egy konkrét \texttt{.bio} programban;
(3) explicit fenyegetési modell, hat támadási forgatókönyv Z3-mal cáfolva;
(4) fordítói korrektség és biszimulációs tétel;
(5) élő kísérleti validáció: 6/6 Z3 PASS, 7/7 rendszerteszt Linux~6.1-en.

A prototípus élesben fut és nyilvánosan elérhető: \bioURL

\appendix

\section{Reprodukálhatóság}

\begin{lstlisting}[style=shellstyle]
# Helyi Z3-ellenorzes (pip install z3-solver szukseges)
python bio_kernel_emu/tests/verify_isomorphism.py
# Vart eredmeny: 6/6 PASS

# Tokioi szerveren
ssh lszok@34.146.249.102 \
  "python3 /home/lszok/bio_kernel_emu_tests/verify_isomorphism.py"

# Kerneltesztek
ssh lszok@34.146.249.102 \
  "cd /home/lszok/dcc_ebpf/src && sudo bash tests/run_proofs.sh"
# Vart eredmeny: 7/7 PASS
\end{lstlisting}

Interaktív demonstráció: \bioURL \quad Z3 bizonyítéklap: \proofURL

\end{magyar}

\end{document}
